%% kapa: exam=2 round=22 year=2011 period=1 subject=appraisal_practice subject_ko=감정평가실무 seq=2 number=2 points=20
%% source: Q-Net 감정평가사 기출문제 (cst00309) artlSeq=5040389 file=extracted/1008/5040389/2011년 제22회 감정평가사 2차시험 기출문제/1교시.hwp
%% conversion: kordoc 4.24.51 -> tools/gen2.py (auto); review=proofread-sonnet; flags=figure
\begin{longquestion}{【문제 2】}
A감정평가법인에 소속된 감정평가사 김정직은 OO도 OO시 OO동 100번지 일대에 소재하는 OO 1-1 주택재개발 정비사업 구역의 사업시행인가 신청을 위한 정비기반시설의 감정평가의뢰를 받았다. 주어진 자료를 활용하여 관련 법규에 따라 주택재개발 정비사업의 시행으로 인하여 용도폐지되는 기존의 정비기반시설 부지와 새로이 설치하는 정비기반시설 예정 부지를 감정평가하시오(단, 가격의 산정과정과 본건 감정평가에 적용할 비교 표준지의 선정이유는 반드시 기술). (20점) \par
- 평가조건 \par
\medskip\noindent (1) 본 감정평가의 대상 토지는 관련 제 규정에 따라 가장 적정하다고 판단되는 인근 공시지가 표준지를 기준으로 공시기준일로부터 가격시점까지의 지가변동률과 토지의 제반가격형성요인을 종합고려하여 적정가격으로 감정평가할 것 \par
\medskip\noindent (2) 국가 또는 지방자치단체로부터 사업시행자인 당해 정비사업조합에 무상으로 양여되는 국ㆍ공유 정비기반시설 부지는 용도폐지를 전제로 감정평가할 것 \par
\medskip\noindent (3) 사업시행자로부터 사업시행인가권자인 지방자치단체에 무상으로 귀속되는 새로이 설치하는 정비기반시설은 토지의 형질변경 등 그 시설의 설치에 소요되는 비용은 포함하지 않고 현장조사 당시 현재의 현황을 기준으로 감정평가할 것 \par
<자료 1> 사업의 개요 \par
1) 정비구역 현황 \par
{\small\centering
\begin{tabular}{|c|c|c|c|c|}
\hline
\makecell{사업의\newline{}종류} & \makecell{구역의\newline{}명칭} & \makecell{위치} & \makecell{면적(m\textsuperscript{2})} & \makecell{비고} \\
\hline
\makecell{주택재개발\newline{}정비사업} & \makecell{OO1-1주택\newline{}재개발정비\newline{}사업구역} & \makecell{OO시 OO동\newline{}100번지 일원} & \makecell{97,600} & \makecell{제2종\newline{}일반주거지역} \\
\hline
\end{tabular}\par}

2) 토지이용 계획 \par
{\small\centering
\begin{tabular}{|c|c|c|c|c|c|}
\hline
\multicolumn{2}{|c|}{\makecell{구분}} & \makecell{명칭} & \makecell{면적(m\textsuperscript{2})} & \makecell{비율(\%)} & \makecell{비고} \\
\hline
\multicolumn{3}{|c|}{\makecell{합계}} & \makecell{97,600} & \makecell{100.0} &  \\
\hline
\multirow{6}{*}{\makecell{토\newline{}지\newline{}이\newline{}용\newline{}계\newline{}획}} & \multirow{3}{*}{\makecell{정비\newline{}기반\newline{}시설 등}} & \makecell{소계} & \makecell{24,800} & \makecell{25.4} &  \\
\cline{3-6}
 &  & \makecell{도로} & \makecell{6,700} & \makecell{6.9} & \makecell{확장 및 신설} \\
\cline{3-6}
 &  & \makecell{공원} & \makecell{18,100} & \makecell{18.5} & \makecell{신설} \\
\cline{2-6}
 & \multirow{3}{*}{\makecell{획지}} & \makecell{소계} & \makecell{72,800} & \makecell{74.6} &  \\
\cline{3-6}
 &  & \makecell{획지1} & \makecell{70,500} & \makecell{72.2} & \makecell{공동주택 및 부대시설} \\
\cline{3-6}
 &  & \makecell{획지2} & \makecell{2,300} & \makecell{2.4} & \makecell{종교시설} \\
\hline
\end{tabular}\par}

3) 건축 계획 \par
{\small\centering
\begin{tabular}{|c|c|c|}
\hline
\multicolumn{2}{|c|}{\makecell{구분}} & \makecell{내용} \\
\hline
\multirow{2}{*}{\makecell{밀도}} & \makecell{건폐율} & \makecell{18.20\%(공동주택 : 15.60\%, 부대시설 : 2.60\%)} \\
\cline{2-3}
 & \makecell{용적율} & \makecell{229.00\%} \\
\hline
\multicolumn{2}{|c|}{\makecell{규모}} & \makecell{공동주택 22개동 1,550세대(16층 ∼ 22층)} \\
\hline
\end{tabular}\par}

4) 지적개황도(축척 없음) \par
\qfigure{figures/02-1.png}
<자료 2> 기타 참고사항 \par
1) 주택재개발 정비구역 지정 고시일 : 2010. 02. 25 \par
2) 주택재개발 정비조합설립인가 고시일 : 2010. 11. 08 \par
3) 가격시점 : 2011. 11. 30 \par
<자료 3> 평가대상 토지 목록 \par
1) 용도 폐지되는 정비기반시설 부지 \par
{\small\centering
\begin{tabular}{|c|c|c|c|c|c|}
\hline
\makecell{일련\newline{}번호} & \makecell{지번} & \makecell{지목} & \makecell{면적(m\textsuperscript{2})} & \makecell{소유자} & \makecell{비고} \\
\hline
\makecell{1} & \makecell{100-28} & \makecell{공원} & \makecell{3,216.0} & \makecell{국\newline{}(국토해양부)} & \makecell{현황\newline{}도시계획시설 공원} \\
\hline
\makecell{2} & \makecell{100-33} & \makecell{도로} & \makecell{1,303.3} & \makecell{OO시} & \makecell{현황\newline{}도시계획시설 도로} \\
\hline
\end{tabular}\par}

2) 새로이 설치되는 정비기반시설 부지 \par
{\small\centering
\begin{tabular}{|c|c|c|c|c|c|}
\hline
\makecell{일련\newline{}번호} & \makecell{지번} & \makecell{지목} & \makecell{면적(m\textsuperscript{2})} & \makecell{소유자} & \makecell{비고} \\
\hline
\makecell{3} & \makecell{100-2} & \makecell{대\newline{}(상업용)} & \makecell{95.6} & \makecell{박부자} & \makecell{대로 3-1호선 확장\newline{}도로 15\%저촉} \\
\hline
\makecell{4} & \makecell{100-5} & \makecell{대\newline{}(주상용)} & \makecell{91.5} & \makecell{강개발} & \makecell{중로 2-8호선 신설\newline{}도로 62\%저촉} \\
\hline
\makecell{5} & \makecell{100-14} & \makecell{대\newline{}(주거용)} & \makecell{138.7} & \makecell{최토지} & \makecell{근린공원신설\newline{}공원 100\%저촉} \\
\hline
\end{tabular}\par}

3) 위 평가 대상 토지 중 일련번호(1) 100-28번지 공원부지의 용도지역은 제1종 일반주거지역이었으나 정비사업의 시행으로 인하여 제2종 일반주거지역으로 변경(2010. 02. 25, 정비구역 지정고시일) 되었음 \par
<자료 4> 토지가격 산정에 참고할 자료 \par
1) 인근 공시지가 표준지 \par
- 당해 정비사업구역 내 공시지가 표준지 \par
{\small\centering
\setlength{\tabcolsep}{3pt}
\begin{tabular}{|c|c|c|c|c|c|c|c|c|c|}
\hline
\multirow{2}{*}{\makecell{일련\newline{}번호}} & \multirow{2}{*}{\makecell{지번}} & \multirow{2}{*}{\makecell{면적\newline{}(m\textsuperscript{2})}} & \multirow{2}{*}{\makecell{이용\newline{}상황}} & \multirow{2}{*}{\makecell{용도\newline{}지역}} & \multirow{2}{*}{\makecell{도로\newline{}교통}} & \multirow{2}{*}{\makecell{형상\newline{}지세}} & \multicolumn{2}{|c|}{\makecell{공시지가(원/m\textsuperscript{2})}} & \multirow{2}{*}{\makecell{비고}} \\
\cline{8-9}
 &  &  &  &  &  &  & \makecell{2010년} & \makecell{2011년} &  \\
\hline
\makecell{371} & \makecell{OO동\newline{}100-3} & \makecell{541.9} & \makecell{상업용} & \makecell{2종\newline{}일주} & \makecell{광대\newline{}소각} & \makecell{정방형\newline{}평지} & \makecell{2,240,000} & \makecell{2,350,000} & \makecell{도로\newline{}15\%저촉} \\
\hline
\makecell{372} & \makecell{OO동\newline{}100-5} & \makecell{147.5} & \makecell{주상\newline{}용} & \makecell{2종\newline{}일주} & \makecell{소로\newline{}각지} & \makecell{정방형\newline{}평지} & \makecell{1,270,000} & \makecell{1,330,000} & \makecell{도로 62\%저촉} \\
\hline
\makecell{373} & \makecell{OO동\newline{}100-12} & \makecell{153.9} & \makecell{단독\newline{}주택} & \makecell{2종\newline{}일주} & \makecell{세로\newline{}(가)} & \makecell{가장형\newline{}평지} & \makecell{1,130,000} & \makecell{1,180,000} & \makecell{공원 100\%저촉} \\
\hline
\end{tabular}\par}

- 당해 정비사업구역 외 공시지가 표준지 \par
{\small\centering
\begin{tabular}{|c|c|c|c|c|c|c|c|c|c|}
\hline
\multirow{2}{*}{\makecell{일련\newline{}번호}} & \multirow{2}{*}{\makecell{지번}} & \multirow{2}{*}{\makecell{면적\newline{}(m\textsuperscript{2})}} & \multirow{2}{*}{\makecell{이용\newline{}상황}} & \multirow{2}{*}{\makecell{용도\newline{}지역}} & \multirow{2}{*}{\makecell{도로\newline{}교통}} & \multirow{2}{*}{\makecell{형상\newline{}지세}} & \multicolumn{2}{|c|}{\makecell{공시지가(원/m\textsuperscript{2})}} & \multirow{2}{*}{\makecell{비고}} \\
\cline{8-9}
 &  &  &  &  &  &  & \makecell{2010년} & \makecell{2011년} &  \\
\hline
\makecell{375} & \makecell{OO동\newline{}110-16} & \makecell{127.9} & \makecell{주상용} & \makecell{2종\newline{}일주} & \makecell{소로\newline{}한면} & \makecell{가장형\newline{}평지} & \makecell{1,230,000} & \makecell{1,290,000} &  \\
\hline
\makecell{377} & \makecell{OO동\newline{}111-5} & \makecell{137.4} & \makecell{단독\newline{}주택} & \makecell{2종\newline{}일주} & \makecell{소로\newline{}한면} & \makecell{세장형\newline{}평지} & \makecell{1,120,000} & \makecell{1,170,000} &  \\
\hline
\makecell{378} & \makecell{OO동\newline{}111-11} & \makecell{600.3} & \makecell{상업\newline{}용} & \makecell{2종\newline{}일주} & \makecell{광대\newline{}세각} & \makecell{세장형\newline{}평지} & \makecell{2,800,000} & \makecell{2,900,000} &  \\
\hline
\makecell{380} & \makecell{OO동\newline{}115-10} & \makecell{109.6} & \makecell{단독\newline{}주택} & \makecell{1종\newline{}일주} & \makecell{세로\newline{}(불)} & \makecell{세장형\newline{}평지} & \makecell{950,000} & \makecell{990,000} &  \\
\hline
\end{tabular}\par}

2) 지가변동률(OO시) \par
{\small\centering
\begin{tabular}{|c|c|c|c|c|}
\hline
\makecell{구분} & \makecell{평균} & \makecell{주거지역} & \makecell{상업지역} & \makecell{공업지역} \\
\hline
\makecell{2010.01.01\newline{}∼ 2010.12.31} & \makecell{1.04583} & \makecell{1.04629} & \makecell{1.04475} & \makecell{1.04400} \\
\hline
\makecell{2011.01.01\newline{}∼ 2011.08.31} & \makecell{1.01591} & \makecell{1.01584} & \makecell{1.01547} & \makecell{1.01622} \\
\hline
\makecell{2011.08.01\newline{}∼ 2011.08.31} & \makecell{1.00173} & \makecell{1.00160} & \makecell{1.00118} & \makecell{1.00191} \\
\hline
\end{tabular}\par}

3) 개별요인 \par
- 도로접면 \par
{\small\centering
\begin{tabular}{|c|c|c|c|c|c|}
\hline
\makecell{구분} & \makecell{광대한면} & \makecell{중로한면} & \makecell{소로한면} & \makecell{세로(가)} & \makecell{세로(불)} \\
\hline
\makecell{격차율} & \makecell{1.00} & \makecell{0.92} & \makecell{0.85} & \makecell{0.80} & \makecell{0.75} \\
\hline
\end{tabular}\par}

※ 도로접면이 각지인 경우는 한면에 접하는 경우에 비해 3\% 우세함 \par
- 이용상황 \par
{\small\centering
\begin{tabular}{|c|c|c|c|}
\hline
\makecell{구분} & \makecell{주거용} & \makecell{주상용} & \makecell{상업용} \\
\hline
\makecell{격차율} & \makecell{1.00} & \makecell{1.08} & \makecell{1.15} \\
\hline
\end{tabular}\par}

- 형상 \par
{\small\centering
\begin{tabular}{|c|c|c|c|c|c|}
\hline
\makecell{구분} & \makecell{가장형} & \makecell{정방형} & \makecell{세장형형} & \makecell{사다리형} & \makecell{부정형} \\
\hline
\makecell{격차율} & \makecell{1.00} & \makecell{0.98} & \makecell{0.96} & \makecell{0.92} & \makecell{0.78} \\
\hline
\end{tabular}\par}

- 도시계획시설 저촉 \par
{\small\centering
\begin{tabular}{|c|c|c|c|}
\hline
\makecell{구분} & \makecell{일반} & \makecell{도로} & \makecell{공원} \\
\hline
\makecell{격차율} & \makecell{1.00} & \makecell{0.85} & \makecell{0.60} \\
\hline
\end{tabular}\par}

\end{longquestion}
